\documentclass[11pt]{article}
\usepackage[margin=1in]{geometry}
\title{ES-FA Code Notes}
\author{Engineer-to-Engineer Notes}
\date{\today}
\begin{document}
\maketitle

\section{paper1\_training/training\_core.py}
\textbf{What}: shared training loop for baseline, experiments, and iterations.\\
\textbf{Why}: keeps software comparisons fair by enforcing one execution path.\\
\textbf{ES-FA fit}: ties Paper 1 training to Paper 2 estimators each batch.\\
\textbf{Interfaces}: \texttt{load\_config()}, \texttt{train\_from\_config()}.\\
\textbf{Tradeoffs}: explicit metrics bookkeeping increases verbosity but improves reproducibility and downstream hardware matching.

\section{paper2\_hardware\_model/estimators.py}
\textbf{What}: proxy estimators for spike cost, memory access, energy, latency.\\
\textbf{Why}: guides hardware-aware training before full FPGA measurement.\\
\textbf{ES-FA fit}: closes software-hardware co-design loop.\\
\textbf{Interfaces}: \texttt{estimate\_spike\_cost}, \texttt{estimate\_memory\_access}, \texttt{estimate\_energy\_proxy}, \texttt{estimate\_latency\_proxy}.\\
\textbf{Tradeoffs}: proxy model is lightweight and interpretable but can underfit hardware non-idealities.

\section{hardware/ and hardware\_validation/kv260}
\textbf{What}: RTL datapath + scheduler blocks and new local validation stack.\\
\textbf{Why}: preserve working RTL while adding research-grade reporting and reproducible build flow.\\
\textbf{ES-FA fit}: enables dense/event comparisons and estimator-vs-hardware checks.\\
\textbf{Interfaces}: top module \texttt{snn\_top}; scripts \texttt{run\_xsim\_regression.py}, \texttt{run\_vivado\_build.py}.\\
\textbf{Tradeoffs}: wrappers avoid risky rewrites, but duplicated orchestration logic must be maintained carefully.

\section{system\_layer and deployment}
\textbf{What}: intent engine, task router, SNN/math/control modules, adaptive policy, CLI.\\
\textbf{Why}: provides a demonstrable system-level contribution beyond isolated kernels.\\
\textbf{ES-FA fit}: adds adaptive execution with explicit policy logs.\\
\textbf{Interfaces}: \texttt{TaskRouter.route()}, \texttt{AdaptiveExecutionPolicy.choose\_mode()}, CLI \texttt{main.py}.\\
\textbf{Tradeoffs}: rule-based intent is easy to debug, but less flexible than learned intent classification.

\section{analysis/evidence\_report.py}
\textbf{What}: produces mandatory evidence tables with absolute/optimized/percent deltas.\\
\textbf{Why}: enforces measurable claims with standardized reporting format.\\
\textbf{ES-FA fit}: codifies four required baseline comparisons.\\
\textbf{Interfaces}: \texttt{generate(results\_root)}.\\
\textbf{Tradeoffs}: strict table schema improves publication readiness but requires clean run metadata.

\section{paper\_output/final\_paper.tex}
\textbf{What}: NeurIPS-style full paper scaffold tied to generated evidence artifacts.\\
\textbf{Why}: avoids drift between implementation and manuscript claims.\\
\textbf{ES-FA fit}: captures method, architecture, experiments, hardware validation, and limitations in one narrative.\\
\textbf{Tradeoffs}: manuscript quality depends on continual sync with latest run outputs.

\end{document}

